\scanpage{165}
\mspage{20}
de la ``formule d'induction'', pour un morphisme $i:s\unsure{\relbar\joinrel\dashv|\longrightarrow}S$, en \deleted{\unsure{formule|formel}}formalisme de dualité :
\[
 i^!(D_S\unsure{P|F})=D_s(i^*(F))
\]
(on prend ici l'inclusion de $s=\operatorname{Spec}k$ dans $S=\operatorname{Spec}V$). Donc le travail à faire concerne les coefficients $p$-primaires, $p=\operatorname{car}\unsure{h|k}>0$. Le cas le plus subtil est celui d'inégales caractéristiques.

b) Le foncteur $R\underline{\Gamma}$ s'obtient en composant $Rj_*$ (où $j:U=\operatorname{Spec}K\mathrel{\relbar\joinrel\dashv}\operatorname{Spec}V=S$ est l'inclusion), avec une sorte de ``foncteur de Greenberg'', version cohomologique.

c) Dans (D) et (E), je me suis \unsure{borné|borner} à parler des coefficients \uline{finis} $F$ -- c'est pour ceux-là que je suis sûr de mon fait. Mais il est certainement vrai que la formule de dualité est encore plus riche, qu'on peut encore énoncer quelque chose \inserted{p.ex} pour \unsure{.} $F$ un schéma en groupes pas néc. fini, p.ex. un schéma abélien \inserted{(avec quelques fibres dégénérées dans (D) ?)} -- mais je n'ai jamais tiré cette question entièrement au clair, même \unsure{heuristiquement}. J'ai une vague souvenance d'un genre de formule
% archive sheet number: 165
